% --- Archon physics formalization source begin ---
% archon:physics
% archon:covers PhyXMiniProblems/problem_phyx_mini_0913.lean
% archon:source-report reports/phyx_mini/problem_phyx_mini_0913.source.json

\chapter{Physics problem phyx\_mini\_0913}
\label{ch:PhyXMiniProblems_problem_phyx_mini_0913}

\paragraph{Problem source.}
\#\# Question

What is the strength of the electric field at the position indicated by the dot?

\#\# Answer choices

- A: A: 1.3 \textbackslash{}times 10\textasciicircum{}\{3\}N/C

- B: B: 5.3 \textbackslash{}times 10\textasciicircum{}\{3\}N/C

- C: C: 7.6 \textbackslash{}times 10\textasciicircum{}\{3\}N/C

- D: D: 3.5 \textbackslash{}times 10\textasciicircum{}\{3\}N/C

\#\# Figure caption (auxiliary; use the image as primary evidence)

The image depicts an arrangement involving electric charges and distances. Here's a detailed description:

- There are two red circles, each containing a plus sign, indicating positive charges of \textbackslash{}(3.0 \textbackslash{}, \textbackslash{}text\{nC\}\textbackslash{}) (nanocoulombs).
- The charges are positioned vertically, with one at the top and the other at the bottom.
- The vertical distance between the two charges is labeled as \textbackslash{}(5.0 \textbackslash{}, \textbackslash{}text\{cm\}\textbackslash{}).
- To the left of the bottom charge, there is a black dot representing a point or possibly a location in space.
- A horizontal arrow points from the black dot toward the left, indicating a distance of \textbackslash{}(5.0 \textbackslash{}, \textbackslash{}text\{cm\}\textbackslash{}) to the left of the top charge.

Overall, the figure represents two identical positive charges placed at a vertical distance with a certain measured point or location horizontally separated from one of them.

\#\# Dataset metadata

Electrostatics; Physical Model Grounding Reasoning, Multi-Formula Reasoning

\paragraph{Recorded answer/context.}
C: C: 7.6 \textbackslash{}times 10\textasciicircum{}\{3\}N/C

\paragraph{Figure/image path.}
/root/proposal\_for\_physic/science-mango/phyx\_mini\_run/phyx\_data/test\_image/913.png

\paragraph{Formalization target.}
create a compiling Lean file with sorry bodies at `PhyXMiniProblems/problem\_phyx\_mini\_0913.lean`. The Lean declarations must preserve the physical quantities, dimensions or dimensional roles, figure labels, governing-law hypotheses, and final relation expressed by this problem.
Use Mathlib/Physlib names found through LeanExplore where available. If a physics API is missing, introduce faithful local abstractions rather than scalar placeholder aliases.

\begin{theorem}[Physics formalization target]
\label{thm:physics:phyx_mini_0913:target}
\lean{PhyXMiniProblems.ProblemPhyXMini0913.problem_phyx_mini_0913}
\uses{def:physics:phyx-mini-0913:phyxminiproblems-problemphyxmini0913-fieldmagnitudeinnewtonspercoulomb, def:physics:phyx-mini-0913:phyxminiproblems-problemphyxmini0913-symmetrictwochargesetup, def:physics:phyx-mini-0913:phyxminiproblems-problemphyxmini0913-matchessuppliedsymmetricchargefigure, def:physics:phyx-mini-0913:phyxminiproblems-problemphyxmini0913-usesvacuumcoulombconstant, def:physics:phyx-mini-0913:phyxminiproblems-problemphyxmini0913-hasphysicaltwochargeparameters, def:physics:phyx-mini-0913:phyxminiproblems-problemphyxmini0913-satisfiespointchargecoulombfieldlaw, def:physics:phyx-mini-0913:phyxminiproblems-problemphyxmini0913-satisfieselectricfieldsuperposition, def:physics:phyx-mini-0913:phyxminiproblems-problemphyxmini0913-resultantmagnituderepresentsfieldatdot, def:physics:phyx-mini-0913:phyxminiproblems-problemphyxmini0913-recordeddatasetanswer, def:physics:phyx-mini-0913:phyxminiproblems-problemphyxmini0913-roundstonearesthundred, def:physics:phyx-mini-0913:phyxminiproblems-problemphyxmini0913-isuniquenearestdisplayedfieldstrength}
The assigned autoformalize agent should translate this physics problem into Lean declarations in the covered file, with theorem and lemma proof bodies written as `by sorry`.
\end{theorem}
\begin{proof}
This is an autoformalization task, not a proof task. Produce statements that can later be proved by the physics prover without weakening the physical model.
\end{proof}

% --- Archon declaration topology begin ---
\section{Declaration topology}
\begin{definition}[Length Quantity]
  \label{def:physics:phyx-mini-0913:phyxminiproblems-problemphyxmini0913-lengthquantity}
  \lean{PhyXMiniProblems.ProblemPhyXMini0913.LengthQuantity}
  A nonnegative, unit-independent physical length
\end{definition}
\begin{proof}
  This is the declared model definition of Length Quantity.
\end{proof}
\begin{definition}[Signed Charge Quantity]
  \label{def:physics:phyx-mini-0913:phyxminiproblems-problemphyxmini0913-signedchargequantity}
  \lean{PhyXMiniProblems.ProblemPhyXMini0913.SignedChargeQuantity}
  A signed, unit-independent physical electric charge
\end{definition}
\begin{proof}
  This is the declared model definition of Signed Charge Quantity.
\end{proof}
\begin{definition}[electric Field Strength Dimension]
  \label{def:physics:phyx-mini-0913:phyxminiproblems-problemphyxmini0913-electricfieldstrengthdimension}
  \lean{PhyXMiniProblems.ProblemPhyXMini0913.electricFieldStrengthDimension}
  The physical dimension M L T⁻² C⁻¹ of electric-field strength
\end{definition}
\begin{proof}
  This is the declared model definition of electric Field Strength Dimension.
\end{proof}
\begin{definition}[Electric Field Magnitude Quantity]
  \label{def:physics:phyx-mini-0913:phyxminiproblems-problemphyxmini0913-electricfieldmagnitudequantity}
  \lean{PhyXMiniProblems.ProblemPhyXMini0913.ElectricFieldMagnitudeQuantity}
  \uses{def:physics:phyx-mini-0913:phyxminiproblems-problemphyxmini0913-electricfieldstrengthdimension}
  A nonnegative, unit-independent electric-field magnitude
\end{definition}
\begin{proof}
  This is the declared model definition of Electric Field Magnitude Quantity.
\end{proof}
\begin{definition}[length In Meters]
  \label{def:physics:phyx-mini-0913:phyxminiproblems-problemphyxmini0913-lengthinmeters}
  \lean{PhyXMiniProblems.ProblemPhyXMini0913.lengthInMeters}
  \uses{def:physics:phyx-mini-0913:phyxminiproblems-problemphyxmini0913-lengthquantity}
  Coherent-SI metre readout of a physical length
\end{definition}
\begin{proof}
  This is the declared model definition of length In Meters.
\end{proof}
\begin{definition}[length In Centimeters]
  \label{def:physics:phyx-mini-0913:phyxminiproblems-problemphyxmini0913-lengthincentimeters}
  \lean{PhyXMiniProblems.ProblemPhyXMini0913.lengthInCentimeters}
  \uses{def:physics:phyx-mini-0913:phyxminiproblems-problemphyxmini0913-lengthquantity, def:physics:phyx-mini-0913:phyxminiproblems-problemphyxmini0913-lengthinmeters}
  Centimetre readout used by the dimension labels in the image
\end{definition}
\begin{proof}
  This is the declared model definition of length In Centimeters.
\end{proof}
\begin{definition}[charge In Coulombs]
  \label{def:physics:phyx-mini-0913:phyxminiproblems-problemphyxmini0913-chargeincoulombs}
  \lean{PhyXMiniProblems.ProblemPhyXMini0913.chargeInCoulombs}
  \uses{def:physics:phyx-mini-0913:phyxminiproblems-problemphyxmini0913-signedchargequantity}
  Coherent-SI coulomb readout of a signed physical charge
\end{definition}
\begin{proof}
  This is the declared model definition of charge In Coulombs.
\end{proof}
\begin{definition}[charge In Nanocoulombs]
  \label{def:physics:phyx-mini-0913:phyxminiproblems-problemphyxmini0913-chargeinnanocoulombs}
  \lean{PhyXMiniProblems.ProblemPhyXMini0913.chargeInNanocoulombs}
  \uses{def:physics:phyx-mini-0913:phyxminiproblems-problemphyxmini0913-signedchargequantity, def:physics:phyx-mini-0913:phyxminiproblems-problemphyxmini0913-chargeincoulombs}
  Nanocoulomb readout used by the two printed charge labels
\end{definition}
\begin{proof}
  This is the declared model definition of charge In Nanocoulombs.
\end{proof}
\begin{definition}[field Magnitude In Newtons Per Coulomb]
  \label{def:physics:phyx-mini-0913:phyxminiproblems-problemphyxmini0913-fieldmagnitudeinnewtonspercoulomb}
  \lean{PhyXMiniProblems.ProblemPhyXMini0913.fieldMagnitudeInNewtonsPerCoulomb}
  \uses{def:physics:phyx-mini-0913:phyxminiproblems-problemphyxmini0913-electricfieldmagnitudequantity}
  Coherent-SI readout of a field magnitude in newtons per coulomb
\end{definition}
\begin{proof}
  This is the declared model definition of field Magnitude In Newtons Per Coulomb.
\end{proof}
\begin{definition}[Source Charge]
  \label{def:physics:phyx-mini-0913:phyxminiproblems-problemphyxmini0913-sourcecharge}
  \lean{PhyXMiniProblems.ProblemPhyXMini0913.SourceCharge}
  The two equal point charges, named by their locations in the image
\end{definition}
\begin{proof}
  This is the declared model definition of Source Charge.
\end{proof}
\begin{definition}[Figure Point]
  \label{def:physics:phyx-mini-0913:phyxminiproblems-problemphyxmini0913-figurepoint}
  \lean{PhyXMiniProblems.ProblemPhyXMini0913.FigurePoint}
  The four geometrically distinguished points in the supplied image
\end{definition}
\begin{proof}
  This is the declared model definition of Figure Point.
\end{proof}
\begin{definition}[expected Source Center]
  \label{def:physics:phyx-mini-0913:phyxminiproblems-problemphyxmini0913-expectedsourcecenter}
  \lean{PhyXMiniProblems.ProblemPhyXMini0913.expectedSourceCenter}
  \uses{def:physics:phyx-mini-0913:phyxminiproblems-problemphyxmini0913-sourcecharge, def:physics:phyx-mini-0913:phyxminiproblems-problemphyxmini0913-figurepoint}
  The center occupied by each named source
\end{definition}
\begin{proof}
  This is the declared model definition of expected Source Center.
\end{proof}
\begin{definition}[Figure Mark Color]
  \label{def:physics:phyx-mini-0913:phyxminiproblems-problemphyxmini0913-figuremarkcolor}
  \lean{PhyXMiniProblems.ProblemPhyXMini0913.FigureMarkColor}
  Colours of the charge circles and observation dot in the raster
\end{definition}
\begin{proof}
  This is the declared model definition of Figure Mark Color.
\end{proof}
\begin{definition}[Figure Charge Sign]
  \label{def:physics:phyx-mini-0913:phyxminiproblems-problemphyxmini0913-figurechargesign}
  \lean{PhyXMiniProblems.ProblemPhyXMini0913.FigureChargeSign}
  The sign glyph drawn in each charge circle
\end{definition}
\begin{proof}
  This is the declared model definition of Figure Charge Sign.
\end{proof}
\begin{definition}[Figure Dimension Segment]
  \label{def:physics:phyx-mini-0913:phyxminiproblems-problemphyxmini0913-figuredimensionsegment}
  \lean{PhyXMiniProblems.ProblemPhyXMini0913.FigureDimensionSegment}
  The three dimension-arrow segments, all labelled 5.0 cm
\end{definition}
\begin{proof}
  This is the declared model definition of Figure Dimension Segment.
\end{proof}
\begin{definition}[Two Charge Field Figure]
  \label{def:physics:phyx-mini-0913:phyxminiproblems-problemphyxmini0913-twochargefieldfigure}
  \lean{PhyXMiniProblems.ProblemPhyXMini0913.TwoChargeFieldFigure}
  \uses{def:physics:phyx-mini-0913:phyxminiproblems-problemphyxmini0913-sourcecharge, def:physics:phyx-mini-0913:phyxminiproblems-problemphyxmini0913-figurepoint, def:physics:phyx-mini-0913:phyxminiproblems-problemphyxmini0913-figuremarkcolor, def:physics:phyx-mini-0913:phyxminiproblems-problemphyxmini0913-figurechargesign, def:physics:phyx-mini-0913:phyxminiproblems-problemphyxmini0913-figuredimensionsegment}
  Literal presentation data transcribed from image 913.png
\end{definition}
\begin{proof}
  This is the declared model definition of Two Charge Field Figure.
\end{proof}
\begin{definition}[position Vector In Meters]
  \label{def:physics:phyx-mini-0913:phyxminiproblems-problemphyxmini0913-positionvectorinmeters}
  \lean{PhyXMiniProblems.ProblemPhyXMini0913.positionVectorInMeters}
  Turn a planar point into its explicitly metre-valued position vector
\end{definition}
\begin{proof}
  This is the declared model definition of position Vector In Meters.
\end{proof}
\begin{definition}[Symmetric Two Charge Setup]
  \label{def:physics:phyx-mini-0913:phyxminiproblems-problemphyxmini0913-symmetrictwochargesetup}
  \lean{PhyXMiniProblems.ProblemPhyXMini0913.SymmetricTwoChargeSetup}
  \uses{def:physics:phyx-mini-0913:phyxminiproblems-problemphyxmini0913-lengthquantity, def:physics:phyx-mini-0913:phyxminiproblems-problemphyxmini0913-signedchargequantity, def:physics:phyx-mini-0913:phyxminiproblems-problemphyxmini0913-electricfieldmagnitudequantity, def:physics:phyx-mini-0913:phyxminiproblems-problemphyxmini0913-sourcecharge, def:physics:phyx-mini-0913:phyxminiproblems-problemphyxmini0913-figurepoint, def:physics:phyx-mini-0913:phyxminiproblems-problemphyxmini0913-twochargefieldfigure}
  Independent physical quantities and fields in the electrostatic setup. The resultant magnitude is an unknown observable. It is related to the norm of the independently modeled resultant vector field only by a separate measurement premise below
\end{definition}
\begin{proof}
  This is the declared model definition of Symmetric Two Charge Setup.
\end{proof}
\begin{definition}[source Position]
  \label{def:physics:phyx-mini-0913:phyxminiproblems-problemphyxmini0913-sourceposition}
  \lean{PhyXMiniProblems.ProblemPhyXMini0913.sourcePosition}
  \uses{def:physics:phyx-mini-0913:phyxminiproblems-problemphyxmini0913-sourcecharge, def:physics:phyx-mini-0913:phyxminiproblems-problemphyxmini0913-symmetrictwochargesetup}
  Spatial position of a named point charge
\end{definition}
\begin{proof}
  This is the declared model definition of source Position.
\end{proof}
\begin{definition}[observation Position]
  \label{def:physics:phyx-mini-0913:phyxminiproblems-problemphyxmini0913-observationposition}
  \lean{PhyXMiniProblems.ProblemPhyXMini0913.observationPosition}
  \uses{def:physics:phyx-mini-0913:phyxminiproblems-problemphyxmini0913-symmetrictwochargesetup}
  Spatial position marked by the black dot
\end{definition}
\begin{proof}
  This is the declared model definition of observation Position.
\end{proof}
\begin{definition}[displacement From Source In Meters]
  \label{def:physics:phyx-mini-0913:phyxminiproblems-problemphyxmini0913-displacementfromsourceinmeters}
  \lean{PhyXMiniProblems.ProblemPhyXMini0913.displacementFromSourceInMeters}
  \uses{def:physics:phyx-mini-0913:phyxminiproblems-problemphyxmini0913-sourcecharge, def:physics:phyx-mini-0913:phyxminiproblems-problemphyxmini0913-positionvectorinmeters, def:physics:phyx-mini-0913:phyxminiproblems-problemphyxmini0913-symmetrictwochargesetup, def:physics:phyx-mini-0913:phyxminiproblems-problemphyxmini0913-sourceposition}
  Displacement from a source to a field point, in metres
\end{definition}
\begin{proof}
  This is the declared model definition of displacement From Source In Meters.
\end{proof}
\begin{definition}[resultant Field At Dot In Newtons Per Coulomb]
  \label{def:physics:phyx-mini-0913:phyxminiproblems-problemphyxmini0913-resultantfieldatdotinnewtonspercoulomb}
  \lean{PhyXMiniProblems.ProblemPhyXMini0913.resultantFieldAtDotInNewtonsPerCoulomb}
  \uses{def:physics:phyx-mini-0913:phyxminiproblems-problemphyxmini0913-symmetrictwochargesetup, def:physics:phyx-mini-0913:phyxminiproblems-problemphyxmini0913-observationposition}
  Resultant field vector at the observation dot, in newtons per coulomb
\end{definition}
\begin{proof}
  This is the declared model definition of resultant Field At Dot In Newtons Per Coulomb.
\end{proof}
\begin{definition}[resultant Field Norm At Dot In Newtons Per Coulomb]
  \label{def:physics:phyx-mini-0913:phyxminiproblems-problemphyxmini0913-resultantfieldnormatdotinnewtonspercoulomb}
  \lean{PhyXMiniProblems.ProblemPhyXMini0913.resultantFieldNormAtDotInNewtonsPerCoulomb}
  \uses{def:physics:phyx-mini-0913:phyxminiproblems-problemphyxmini0913-symmetrictwochargesetup, def:physics:phyx-mini-0913:phyxminiproblems-problemphyxmini0913-resultantfieldatdotinnewtonspercoulomb}
  Norm of the resultant field vector at the dot, in newtons per coulomb
\end{definition}
\begin{proof}
  This is the declared model definition of resultant Field Norm At Dot In Newtons Per Coulomb.
\end{proof}
\begin{definition}[Matches Supplied Symmetric Charge Figure]
  \label{def:physics:phyx-mini-0913:phyxminiproblems-problemphyxmini0913-matchessuppliedsymmetricchargefigure}
  \lean{PhyXMiniProblems.ProblemPhyXMini0913.MatchesSuppliedSymmetricChargeFigure}
  \uses{def:physics:phyx-mini-0913:phyxminiproblems-problemphyxmini0913-lengthinmeters, def:physics:phyx-mini-0913:phyxminiproblems-problemphyxmini0913-lengthincentimeters, def:physics:phyx-mini-0913:phyxminiproblems-problemphyxmini0913-chargeinnanocoulombs, def:physics:phyx-mini-0913:phyxminiproblems-problemphyxmini0913-expectedsourcecenter, def:physics:phyx-mini-0913:phyxminiproblems-problemphyxmini0913-positionvectorinmeters, def:physics:phyx-mini-0913:phyxminiproblems-problemphyxmini0913-symmetrictwochargesetup}
  Literal image evidence and its association with the physical quantities and planar coordinates. No resultant field value occurs in this predicate
\end{definition}
\begin{proof}
  This is the declared model definition of Matches Supplied Symmetric Charge Figure.
\end{proof}
\begin{definition}[Uses Vacuum Coulomb Constant]
  \label{def:physics:phyx-mini-0913:phyxminiproblems-problemphyxmini0913-usesvacuumcoulombconstant}
  \lean{PhyXMiniProblems.ProblemPhyXMini0913.UsesVacuumCoulombConstant}
  \uses{def:physics:phyx-mini-0913:phyxminiproblems-problemphyxmini0913-symmetrictwochargesetup}
  Coherent-SI calibration of Physlib's Coulomb constant. The scalar has units N m²/C² in the field law below
\end{definition}
\begin{proof}
  This is the declared model definition of Uses Vacuum Coulomb Constant.
\end{proof}
\begin{definition}[Has Physical Two Charge Parameters]
  \label{def:physics:phyx-mini-0913:phyxminiproblems-problemphyxmini0913-hasphysicaltwochargeparameters}
  \lean{PhyXMiniProblems.ProblemPhyXMini0913.HasPhysicalTwoChargeParameters}
  \uses{def:physics:phyx-mini-0913:phyxminiproblems-problemphyxmini0913-lengthinmeters, def:physics:phyx-mini-0913:phyxminiproblems-problemphyxmini0913-symmetrictwochargesetup, def:physics:phyx-mini-0913:phyxminiproblems-problemphyxmini0913-observationposition, def:physics:phyx-mini-0913:phyxminiproblems-problemphyxmini0913-displacementfromsourceinmeters}
  Positivity and source-observation separation conditions
\end{definition}
\begin{proof}
  This is the declared model definition of Has Physical Two Charge Parameters.
\end{proof}
\begin{definition}[Satisfies Point Charge Coulomb Field Law]
  \label{def:physics:phyx-mini-0913:phyxminiproblems-problemphyxmini0913-satisfiespointchargecoulombfieldlaw}
  \lean{PhyXMiniProblems.ProblemPhyXMini0913.SatisfiesPointChargeCoulombFieldLaw}
  \uses{def:physics:phyx-mini-0913:phyxminiproblems-problemphyxmini0913-chargeincoulombs, def:physics:phyx-mini-0913:phyxminiproblems-problemphyxmini0913-symmetrictwochargesetup, def:physics:phyx-mini-0913:phyxminiproblems-problemphyxmini0913-sourceposition, def:physics:phyx-mini-0913:phyxminiproblems-problemphyxmini0913-displacementfromsourceinmeters}
  For charge q and displacement r from the source to a field point, the electric field is k q r / ‖r‖³. This law is quantified over every point away from each source and contains no requested numerical result
\end{definition}
\begin{proof}
  This is the declared model definition of Satisfies Point Charge Coulomb Field Law.
\end{proof}
\begin{definition}[Satisfies Electric Field Superposition]
  \label{def:physics:phyx-mini-0913:phyxminiproblems-problemphyxmini0913-satisfieselectricfieldsuperposition}
  \lean{PhyXMiniProblems.ProblemPhyXMini0913.SatisfiesElectricFieldSuperposition}
  \uses{def:physics:phyx-mini-0913:phyxminiproblems-problemphyxmini0913-sourcecharge, def:physics:phyx-mini-0913:phyxminiproblems-problemphyxmini0913-symmetrictwochargesetup}
  The resultant field is the vector sum of the two source fields
\end{definition}
\begin{proof}
  This is the declared model definition of Satisfies Electric Field Superposition.
\end{proof}
\begin{definition}[Resultant Magnitude Represents Field At Dot]
  \label{def:physics:phyx-mini-0913:phyxminiproblems-problemphyxmini0913-resultantmagnituderepresentsfieldatdot}
  \lean{PhyXMiniProblems.ProblemPhyXMini0913.ResultantMagnitudeRepresentsFieldAtDot}
  \uses{def:physics:phyx-mini-0913:phyxminiproblems-problemphyxmini0913-fieldmagnitudeinnewtonspercoulomb, def:physics:phyx-mini-0913:phyxminiproblems-problemphyxmini0913-symmetrictwochargesetup, def:physics:phyx-mini-0913:phyxminiproblems-problemphyxmini0913-resultantfieldnormatdotinnewtonspercoulomb}
  The dimensionful scalar observable is the Euclidean magnitude of the modeled resultant vector at the black dot. This is a general measurement relation, not a numerical answer
\end{definition}
\begin{proof}
  This is the declared model definition of Resultant Magnitude Represents Field At Dot.
\end{proof}
\begin{definition}[Answer Choice]
  \label{def:physics:phyx-mini-0913:phyxminiproblems-problemphyxmini0913-answerchoice}
  \lean{PhyXMiniProblems.ProblemPhyXMini0913.AnswerChoice}
  Labels attached to the four multiple-choice answers
\end{definition}
\begin{proof}
  This is the declared model definition of Answer Choice.
\end{proof}
\begin{definition}[field Strength In Newtons Per Coulomb]
  \label{def:physics:phyx-mini-0913:phyxminiproblems-problemphyxmini0913-answerchoice-fieldstrengthinnewtonspercoulomb}
  \lean{PhyXMiniProblems.ProblemPhyXMini0913.AnswerChoice.fieldStrengthInNewtonsPerCoulomb}
  \uses{def:physics:phyx-mini-0913:phyxminiproblems-problemphyxmini0913-answerchoice}
  Field strength in newtons per coulomb printed by each answer choice
\end{definition}
\begin{proof}
  This is the declared model definition of field Strength In Newtons Per Coulomb.
\end{proof}
\begin{definition}[recorded Dataset Answer]
  \label{def:physics:phyx-mini-0913:phyxminiproblems-problemphyxmini0913-recordeddatasetanswer}
  \lean{PhyXMiniProblems.ProblemPhyXMini0913.recordedDatasetAnswer}
  \uses{def:physics:phyx-mini-0913:phyxminiproblems-problemphyxmini0913-answerchoice}
  The answer label recorded by the supplied dataset metadata
\end{definition}
\begin{proof}
  This is the declared model definition of recorded Dataset Answer.
\end{proof}
\begin{definition}[Rounds To Nearest Hundred]
  \label{def:physics:phyx-mini-0913:phyxminiproblems-problemphyxmini0913-roundstonearesthundred}
  \lean{PhyXMiniProblems.ProblemPhyXMini0913.RoundsToNearestHundred}
  A value rounds to a displayed reading at the nearest 100 N/C
\end{definition}
\begin{proof}
  This is the declared model definition of Rounds To Nearest Hundred.
\end{proof}
\begin{definition}[Is Unique Nearest Displayed Field Strength]
  \label{def:physics:phyx-mini-0913:phyxminiproblems-problemphyxmini0913-isuniquenearestdisplayedfieldstrength}
  \lean{PhyXMiniProblems.ProblemPhyXMini0913.IsUniqueNearestDisplayedFieldStrength}
  \uses{def:physics:phyx-mini-0913:phyxminiproblems-problemphyxmini0913-answerchoice}
  A displayed choice is strictly nearer than every alternative
\end{definition}
\begin{proof}
  This is the declared model definition of Is Unique Nearest Displayed Field Strength.
\end{proof}
\begin{lemma}[resultant Field At Dot eq coulomb Sum]
  \label{lem:physics:phyx-mini-0913:phyxminiproblems-problemphyxmini0913-resultantfieldatdot-eq-coulombsum}
  \lean{PhyXMiniProblems.ProblemPhyXMini0913.resultantFieldAtDot_eq_coulombSum}
  \uses{def:physics:phyx-mini-0913:phyxminiproblems-problemphyxmini0913-chargeincoulombs, def:physics:phyx-mini-0913:phyxminiproblems-problemphyxmini0913-sourcecharge, def:physics:phyx-mini-0913:phyxminiproblems-problemphyxmini0913-symmetrictwochargesetup, def:physics:phyx-mini-0913:phyxminiproblems-problemphyxmini0913-observationposition, def:physics:phyx-mini-0913:phyxminiproblems-problemphyxmini0913-displacementfromsourceinmeters, def:physics:phyx-mini-0913:phyxminiproblems-problemphyxmini0913-resultantfieldatdotinnewtonspercoulomb, def:physics:phyx-mini-0913:phyxminiproblems-problemphyxmini0913-hasphysicaltwochargeparameters, def:physics:phyx-mini-0913:phyxminiproblems-problemphyxmini0913-satisfiespointchargecoulombfieldlaw, def:physics:phyx-mini-0913:phyxminiproblems-problemphyxmini0913-satisfieselectricfieldsuperposition}
  Coulomb's vector law and superposition give the source-by-source expression at the observation dot. It contains only source data and the governing constant
\end{lemma}
\begin{proof}
  The conclusion follows from the governing relations and figure readouts named in the statement of resultant Field At Dot eq coulomb Sum.
\end{proof}
\begin{lemma}[resultant Field At Dot vertical Component eq zero]
  \label{lem:physics:phyx-mini-0913:phyxminiproblems-problemphyxmini0913-resultantfieldatdot-verticalcomponent-eq-zero}
  \lean{PhyXMiniProblems.ProblemPhyXMini0913.resultantFieldAtDot_verticalComponent_eq_zero}
  \uses{def:physics:phyx-mini-0913:phyxminiproblems-problemphyxmini0913-symmetrictwochargesetup, def:physics:phyx-mini-0913:phyxminiproblems-problemphyxmini0913-resultantfieldatdotinnewtonspercoulomb, def:physics:phyx-mini-0913:phyxminiproblems-problemphyxmini0913-matchessuppliedsymmetricchargefigure, def:physics:phyx-mini-0913:phyxminiproblems-problemphyxmini0913-hasphysicaltwochargeparameters, def:physics:phyx-mini-0913:phyxminiproblems-problemphyxmini0913-satisfiespointchargecoulombfieldlaw, def:physics:phyx-mini-0913:phyxminiproblems-problemphyxmini0913-satisfieselectricfieldsuperposition}
  By the reflection symmetry of the two equal charges, the vertical components of their fields cancel at the observation dot
\end{lemma}
\begin{proof}
  The conclusion follows from the governing relations and figure readouts named in the statement of resultant Field At Dot vertical Component eq zero.
\end{proof}
\begin{lemma}[resultant Field Norm At Dot numerical Bounds]
  \label{lem:physics:phyx-mini-0913:phyxminiproblems-problemphyxmini0913-resultantfieldnormatdot-numericalbounds}
  \lean{PhyXMiniProblems.ProblemPhyXMini0913.resultantFieldNormAtDot_numericalBounds}
  \uses{def:physics:phyx-mini-0913:phyxminiproblems-problemphyxmini0913-symmetrictwochargesetup, def:physics:phyx-mini-0913:phyxminiproblems-problemphyxmini0913-resultantfieldnormatdotinnewtonspercoulomb, def:physics:phyx-mini-0913:phyxminiproblems-problemphyxmini0913-matchessuppliedsymmetricchargefigure, def:physics:phyx-mini-0913:phyxminiproblems-problemphyxmini0913-usesvacuumcoulombconstant, def:physics:phyx-mini-0913:phyxminiproblems-problemphyxmini0913-hasphysicaltwochargeparameters, def:physics:phyx-mini-0913:phyxminiproblems-problemphyxmini0913-satisfiespointchargecoulombfieldlaw, def:physics:phyx-mini-0913:phyxminiproblems-problemphyxmini0913-satisfieselectricfieldsuperposition}
  The two Coulomb contributions yield an unrounded magnitude of approximately 7.63 * 10\textasciicircum{}3 N/C, which lies in the nearest-hundred rounding interval for 7.6 * 10\textasciicircum{}3 N/C
\end{lemma}
\begin{proof}
  The conclusion follows from the governing relations and figure readouts named in the statement of resultant Field Norm At Dot numerical Bounds.
\end{proof}
% --- Archon declaration topology end ---
% --- Archon physics formalization source end ---
